本作品面向医院侧数据与模型提供方参数同时敏感的协同推理场景，构建了一套可运行、可验证、可复现的双向隐私安全推理系统。针对传统单向隐私保护难以同时兼顾数据合规与模型知识产权保护的问题，本作品以 DynamicViT 为载体，将原始动态词元剪枝中的删除式控制流改写为固定图下可执行的密态语义，通过安全比较、安全选择与稳定排序机制，使动态剪枝能力能够进入两方安全计算链路。

在系统实现上，本作品将浏览器端本地预处理、秘密分享生成、服务端权威快检、SPU 密态前向计算与审计记录整合为端到端闭环。原始影像仅在终端侧处理，跨越信任边界后仅传输密态分片、结构化元数据与审计摘要；服务端在进入密态计算前完成协议、张量与质量摘要检查，从而避免将输入天然可信作为隐含前提。

实验结果表明，本作品在 524 张医疗验证样本上取得 92.7481\% 的阈值精度，较同仓固定结构静态对照线提升 0.7634 个百分点，较未经安全化边界重写的原始 DynamicViT 提升 17.3664 个百分点。端到端双向隐私原型在 32 条医疗部署验证样本上的平均时延为 89.06 秒/样本，双向总通信量为 84.47 GiB。系统还通过协议异常输入、重放与并发守卫等 17 类黑盒验证，为安全边界和原型可服务性提供了证据支撑。
